\section{O Marco Normativo Brasileiro e o Panorama Global}

A consolidação de gêmeos digitais no Brasil encontra-se em um ponto de inflexão regulatório. Para que tecnologias de fronteira operem de forma segura e interoperável em escala nacional, especialmente no âmbito do SUS, é mandatório o estabelecimento de parâmetros técnicos, terminológicos e ontológicos rigorosos. 

Em fevereiro de 2026, a Associação Brasileira de Normas Técnicas (ABNT), impulsionada por debates na comunidade tecnológica e na Associação Brasileira de Internet das Coisas (ABINC), publicou a norma \textbf{ABNT NBR ISO/IEC 3173}~\cite{abinc2026brasil}. Este documento estabeleceu o vocabulário padronizado, os domínios conceituais e as diretrizes arquitetônicas essenciais para o desenvolvimento de gêmeos digitais em solo brasileiro. A norma atua como uma âncora regulatória fundamental, mitigando a fragmentação do mercado e fornecendo a base técnica necessária para futuras regulamentações sanitárias pela ANVISA.

\destaque{A padronização semântica e estrutural promovida pela NBR ISO/IEC 3173 é o pré-requisito técnico para a criação de sistemas integrados capazes de processar os gêmeos digitais odontológicos de milhões de brasileiros sem colapsos de interoperabilidade.}

Sob a perspectiva econômica e estratégica, a adoção destas normas alinha o Brasil às macrotendências globais em \textit{Digital Health}. Projeções financeiras ilustram a magnitude desta transformação: o mercado global de gêmeos digitais em saúde foi avaliado em aproximadamente US\$ 24,48 bilhões em 2025, com uma expectativa de crescimento exponencial que deve atingir a marca de US\$ 384,79 bilhões até 2034, representando uma Taxa de Crescimento Anual Composta (CAGR) de 35,4\%~\cite{fortune2026digital}. 

Este volume de investimentos reflete a percepção do mercado sobre a inevitabilidade da transição para a saúde de precisão preditiva. Para o Brasil, contudo, o desafio não é apenas acompanhar este crescimento, mas direcioná-lo de modo que a infraestrutura desenvolvida atenda primordialmente às necessidades epidemiológicas da população, integrando plataformas de código aberto como o OpenCode Ecosystem~\cite{robles2023opentwins} para evitar o aprisionamento tecnológico (\textit{vendor lock-in}) imposto por corporações multinacionais de softwares proprietários.
